%\cvsection{Motivation}
\begin{cventries}
\cventry
    {} % Empty position
    {Motivation} % Project
    {} % Empty location
    {} % Empty date
    {
      %Guten Tag, \newline
      %hiermit bewerbe ich mich für die Stelle \textit{Studentische Hilfskraft (Programmierer/-in)}.\newline
      Ich bin momentan auf der Zielgeraden meines Masterstudiums und fühlte mich von der flexiblen Position angesprochen, da sich so Studium und Job hervorragend kombinieren lassen und die Vorlesungsfreie Zeit sinnvoll von mir genutzt werden kann.\newline
      Zwar habe ich noch nicht viel berufliche Erfahrung, habe aber bereits einige kleinere Projekte gemacht, die wahrscheinlich ziemlich gut auf die Anforderungen passen. 
      Die meisten davon können auch unter meinem oben verlinkten GitHub-Account nachvollzogen werden.
      Außerdem habe ich mein Praktikum in einem StartUp gemacht, bei dem ich einen Großteil der Zeit mit PHP gearbeitet hab. Wobei wir, um ehrlich zu sein, dort dann auf ein anderes, moderneres Framework (Nuxt.js) gewechselt haben. Dementsprechend sind meine PHP-Kenntnisse etwas eingerostet, lassen sich aber sicherlich schnell wieder auffrischen. \newline
      Eine Migration zu Laravel wäre zwar nicht mein Hauptinteresse, damit habe ich auch noch nie gearbeitet, bin aber durchaus offen für Neues.\newline
      Ein erhöhtes Interesse hätte ich an der Umsetzung von Smartphone-Apps. Hier könnte ich auch einige Erfahrung, sowohl im nativen Bereich (hierbei vor allem Swift(UI), hätte aber auch Interesse Kotlin zu lernen), als auch im Cross-Plattform-Bereich mit Flutter, mitbringen. Eine sehr minimalistische Passwort-Check Web-App habe ich auch als TWA im Playstore \href{https://play.google.com/store/apps/details?id=hannepps.tools.passwordchecker}{veröffentlicht}. \newline
      Vorallem freue ich mich aber auf ein abwechslungsreiches, frohes Arbeitsumfeld!
    }
\end{cventries}